\section{4.3.Runner}\label{runner}

\subsection{4.3. Módulo Runner}\label{muxf3dulo-runner}

El módulo Runner constituye la capa encargada de la interacción con las
herramientas externas utilizadas durante el proceso de análisis. Su
principal responsabilidad consiste en gestionar la ejecución de Slither,
Mythril y Echidna, así como preparar el entorno necesario para
garantizar la correcta realización del análisis.

La existencia de un módulo específico para esta tarea permite desacoplar
la lógica propia de EVMAudit de las particularidades de cada
herramienta, facilitando la incorporación de nuevas soluciones de
análisis en futuras versiones.

!TODO: insertar referencia a la Figura X (diagrama de arquitectura).

\subsubsection{4.3.1. Responsabilidades del
módulo}\label{responsabilidades-del-muxf3dulo}

Las principales responsabilidades implementadas por el módulo Runner son
las siguientes:

\begin{itemize}
\tightlist
\item
  detección automática del nombre del contrato
\item
  configuración de la versión del compilador Solidity
\item
  ejecución de Slither
\item
  ejecución de Mythril
\item
  ejecución de Echidna
\item
  almacenamiento de resultados intermedios
\item
  gestión de errores producidos durante la ejecución.
\end{itemize}

Gracias a esta aproximación, el resto de módulos del sistema pueden
trabajar de forma independiente sin necesidad de conocer los detalles
específicos de cada herramienta externa.

\subsubsection{4.3.2. Detección automática del nombre del
contrato}\label{detecciuxf3n-automuxe1tica-del-nombre-del-contrato}

Durante las primeras fases de desarrollo se observó que el nombre del
fichero Solidity no siempre coincide con el nombre del contrato definido
en su interior.

Por este motivo se implementó un mecanismo de detección automática que
analiza el código fuente y extrae el nombre real del contrato utilizando
expresiones regulares.

Esta decisión permite evitar errores durante las fases posteriores del
análisis y elimina la necesidad de que el usuario proporcione
manualmente dicha información.

Además, este mecanismo resulta especialmente útil cuando la librería es
utilizada desde otras aplicaciones externas, ya que reduce la cantidad
de parámetros que deben proporcionarse.

!TODO: insertar fragmento de código de
\texttt{detect\_contract\_name()}.

\subsubsection{4.3.3. Configuración automática del
compilador}\label{configuraciuxf3n-automuxe1tica-del-compilador}

Las herramientas de análisis empleadas dependen de la versión del
compilador Solidity utilizada por el contrato.

Con el objetivo de garantizar la reproducibilidad del análisis y evitar
incompatibilidades entre versiones, EVMAudit analiza automáticamente la
directiva \texttt{pragma\ solidity} presente en el contrato y configura
la versión correspondiente mediante la utilidad \texttt{solc-select}.

Este mecanismo permite adaptar dinámicamente el entorno de ejecución y
facilita el análisis de contratos desarrollados con diferentes versiones
del lenguaje.

!TODO: insertar referencia cruzada a la Sección 2.X (Solidity y
compilador).

!TODO: insertar fragmento de código de \texttt{\_set\_solc\_version()}.

\subsubsection{4.3.4. Ejecución de
Slither}\label{ejecuciuxf3n-de-slither}

Slither constituye la primera herramienta utilizada durante el proceso
de análisis debido a su rapidez y a la gran cantidad de detectores
disponibles.

La función encargada de su ejecución verifica previamente que la
herramienta se encuentre instalada y posteriormente lanza el análisis
sobre el contrato proporcionado.

Una vez finalizada la ejecución, la salida obtenida se almacena en
formato JSON para garantizar la trazabilidad de los resultados.

Durante la implementación fue necesario contemplar ciertos
comportamientos específicos de Slither. En particular, la herramienta
devuelve el código de salida 255 cuando detecta vulnerabilidades, aunque
dicho comportamiento no representa un error real.

Por este motivo, EVMAudit incorpora una lógica específica para
interpretar correctamente este caso y continuar con el proceso de
análisis.

!TODO: insertar fragmento de código de \texttt{run\_slither()}.

!TODO: referencia cruzada a la Sección 2.X (Slither).

\subsubsection{4.3.5. Ejecución de
Mythril}\label{ejecuciuxf3n-de-mythril}

La segunda fase del análisis se realiza mediante Mythril, herramienta
basada en ejecución simbólica.

A diferencia de Slither, Mythril presenta tiempos de ejecución
superiores y requiere parámetros adicionales relacionados con la
profundidad máxima de exploración y el tiempo máximo permitido para
completar el análisis.

Durante el desarrollo se incorporaron mecanismos para:

\begin{itemize}
\tightlist
\item
  limitar el tiempo máximo de ejecución
\item
  controlar posibles errores durante el análisis
\item
  procesar las salidas generadas por la herramienta
\item
  preservar los resultados originales.
\end{itemize}

Estas medidas permiten evitar que un contrato especialmente complejo
bloquee el proceso completo de análisis.

!TODO: insertar fragmento de código de \texttt{run\_mythril()}.

!TODO: referencia cruzada a la Sección 2.X (Mythril).

\subsubsection{4.3.6. Ejecución de
Echidna}\label{ejecuciuxf3n-de-echidna}

Una vez generadas las propiedades correspondientes, EVMAudit ejecuta
Echidna para realizar una fase adicional de validación mediante fuzzing.

La integración de Echidna presentó una complejidad superior a la
observada en Slither y Mythril debido a ciertas particularidades de la
herramienta y a las diferencias existentes entre versiones.

Durante el desarrollo fue necesario incorporar mecanismos específicos
para:

\begin{itemize}
\tightlist
\item
  reconstruir nombres de propiedades
\item
  interpretar adecuadamente las salidas generadas
\item
  distinguir entre pruebas superadas y vulnerabilidades explotables
\item
  gestionar posibles inconsistencias en los resultados.
\end{itemize}

Estas adaptaciones permitieron integrar Echidna dentro del flujo general
de EVMAudit sin alterar el resto de componentes de la arquitectura.

!TODO: insertar fragmento de código de \texttt{run\_echidna()}.

!TODO: referencia cruzada a la Sección 2.X (Echidna).

\subsubsection{4.3.7. Persistencia de resultados
intermedios}\label{persistencia-de-resultados-intermedios}

Con el fin de favorecer la trazabilidad y la reproducibilidad del
análisis, EVMAudit almacena las salidas originales generadas por las
distintas herramientas utilizadas.

Esta decisión permite:

\begin{itemize}
\tightlist
\item
  inspeccionar manualmente los resultados obtenidos
\item
  reproducir análisis posteriores
\item
  depurar posibles errores
\item
  conservar evidencias originales.
\end{itemize}

La preservación de las salidas originales resulta especialmente
relevante en el contexto de auditorías de seguridad, donde la
trazabilidad de los hallazgos constituye un aspecto fundamental.
